\cotefile{157-4}
\scanpage{18}
\mspage{12}

\unsure{.} Il ne semble pas que les axiomes
$M0$ à $M4$, plus $M'$ et $M''$, suffisent à
développer les suites exactes de suspension et
de cosuspension. Il faut un petit
\unsure{chouïa} de plus :

\uline{Définition 2} J'appelle \uline{pré-$Q$-catégorie de modèles à gauche}, la donnée de \deleted{\unsure{cat.}\ \illegible{1}} $M$ et de
$W,C,F,TC,TF\subset\mathrm{Fl}\,M$, satisfaisant les axiomes
$M0$ à $M5$, $M'$, $M''$, plus la forme renforcée
de la \inserted{première} inclusion de $M'$
\[
 TC=C\cap W\qquad\left(\begin{aligned}&\text{non seulement}\\&TC\subset C\cap W\end{aligned}\right) \tag*{$M'_0(g)$}
\]

De même, on définit la \uline{pré-$Q$-catégorie de modèles à droite}, par la condition
\[
 TF=F\cap W. \tag*{$M'_0(d)$}
\]

\unsure{NB} Quand les deux conditions sont satisfaites
à la fois, \unsure{les|des} \inserted{la donnée de} $TC,TF$ résulte\deleted{\unsure{nt}} déjà de
celle de $W,C,F$, et on a simplement
une \inserted{$Q$-}cat. de Modèles ``propre'' (au sens
de \unsure{Quillen}).

\uline{Scholie} \deleted{Il \unsure{me semble} que} \inserted{\unsure{Pour moi,}} l'introduction
de $\unsure{E|C},TC,\deleted{\unsure{F}}F,TF$ est purement auxiliaire
--- ce qui m'importe, ce sont les propriétés
\unsure{de|des} $W$. Ainsi, l'existence de $C,TC,F,TF$
satisfaisant aux axiomes faibles $M0$ à $M4$
$M',M''$, du fait qu'elle \deleted{\unsure{fait}} donne les factorisations
% archive sheet number: 17
\scanpage{19}
\mspage{13}

dont il est question dans prop 3, \deleted{fait} entraîne
la propriété remarquable de $W$ énoncée
dans cette proposition, et \unsure{dans|dès} le corollaire
qui la précède. \deleted{\unsure{Et} \unsure{=}} De \unsure{$=$}, ces propriétés
de factorisation\textsuperscript{$*$} assurent l'existence d'une
% slanted
\margin{\unsure{sous une forme} légèrement renforcée, il est vrai --- \unsure{i.e.} que ces factorisations \unsure{dép\textsuperscript{dt}} \uline{\unsure{fonctoriellt}} \unsure{p.\ en $f$}\ldots}
bonne théorie des fibres \inserted{et cofibres} $W$-homotopiques,
sans \inserted{besoin de} référence dans cette théorie à $C,TC,F,TF$
(\deleted{qui} dont l'introduction donne \inserted{\unsure{du reste}} des \unsure{petites|petits} \unsure{complications|compléments}\ldots). Mais je ne vois pas comment,
à partir des seules \unsure{hyps|hyp.} de factorisation,
prouver la propriété \unsure{crucial|cruciale} d'exactitude
pour les carrés $W$-cartésiens, et
$W$-cocartésiens, nécessaire pour pouvoir
développer la suite \uuline{exacte} de cosuspension
et de suspension. C'est pour avoir cette
exactitude que l'\unsure{introduction|intr.duction} auxiliaire,
et transitoire, des \unsure{``attirail''|l'attirail|attirails|adjuvants} $(C,TC,F,TF)$
est utile, l'axiome \inserted{suppl.} à gauche $TC=C\cap W$
assurant (sauf erreur) l'exactitude des
carrés \inserted{$W$-}cartésiens, et l'axiome à droite
$TF=F\cap W$ celle des carrés $W$-cocartésiens.
(Chose à vérifier avec plus de soin que je
ne l'ai fait à présent.) Une fois ceci obtenu,
et éventuellement idem pour les $(M(I),W(I))$
pour toute catégorie-diagramme $I$, on peut
% archive sheet number: 18
\scanpage{20}
\mspage{14}

totalement oublier ces ensembles auxiliaires
--- \unsure{j'ai envie|j'aimerais} en tous cas de les oublier.

Mon espoir, c'est que pour une $W$ donnée,
même si on n'arrive pas à trouver
une structure de triple de Quillen
$(W,C,F)$ associée, (p.ex. parce qu'il n'y en
aurait pas), on peut trouver \inserted{à la fois} des
quintuplets \inserted{$(W,C,F,TC,TF)$} qui \unsure{sont|soient} des pré-Quillen à \unsure{g},
et des quintuplets \inserted{(\unsure{d'}\uline{autres} quintuplets \unsure{!|?})} qui \unsure{soient|sont} des pré-Quillen
à droite.

Voici le genre d'énoncé en forme,
pour le moment conjectural voire totalement
heuristique, auquel je songe :

\uline{Théorème-Scholie conjectural} : On suppose
que $M$ est stable par $\varinjlim$ ``quelconques'',
et par $\varprojlim$ finies, et que $W$ est
stable par $\varinjlim$ filtrantes \deleted{\unsure{cela implique}} \inserted{(en plus de la \unsure{rétroversion} faible)}.

\textcircled{1} Alors $\mathrm{Cof}_{W}$ et $\deleted{\unsure{\mathrm{Cof}}}W^{\unsure{\mathrm{univ}}}$ ont les
propriétés de stabilité a) \unsure{et c)|à e)} de la
prop. 2 (p.~4), et aussi la propriété de
stabilité b par facteurs directs, pourvu
que $W$ lui-même l'a.

\textcircled{2} Supposons $M$ \uline{accessible}, de sorte qu'on
dispose du th. 1. \deleted{Pour que} Conditions
équivalentes sur $W$ :
% archive sheet number: 19
\scanpage{21}
\mspage{15}

a) Il existe une partie \unsure{petite|petisse} $C_0$ de $\mathrm{Cof}_W$, et $TC_0$ de
$W^{\mathrm{univ}}$, telles que
\[
 TF\overset{\mathrm{def}}{=}C_{0*}\subset\deleted{\unsure{\mathrm{Fib}_W}}W,\qquad
 (TC_0)_*\overset{\mathrm{def}}{=}F\subset\underline{\mathrm{Fib}}_W
\]

\margin{$W^{\mathrm{univ}}\subset\mathrm{Cof}_W$}
a$'$) Il existe un quadruplet $(C,TC,F,TF)$ de
parties de $\mathrm{Fl}(M)$, satisfaisant les conditions
$M1$ à $M4$, $M',M''$ de la page 6.
\margin{avec $C,TC$ clos à gauche, $F,TF$ clos à droite}

a$''$) \unsure{Idem} a$'$), en supposant $M'$ sous forme
renforcée \unsure{$\nearrow$|à} \unsure{comme|en} $M'_0(g)$, i.e. on a une structure
de pré-$Q$-cat. de modèles à \unsure{g.} $(TC=C\cap W)$.

b) On a $\unsure{C|\subset}$
\[
 (\mathrm{Cof}_W)_*\subset\deleted{\unsure{\mathrm{Fib}_W}}W,\quad
 (W^{\mathrm{univ}})_*\subset\underset{\deleted{\unsure{W}}}{\mathrm{Fib}}_W
\]
(donc $(\mathrm{Cof}_W)_*\subset\mathrm{Fib}_W\cap W$ puisque
$(\mathrm{Cof}_W)_*\subset(W^{\mathrm{univ}})_*$\unsure{).}

b$'$) On a
\[
 (\mathrm{Cof}_W)_*\subset W_{\mathrm{univ}},\quad
 (W^{\mathrm{univ}})_*\subset\mathrm{Fib}_W.
\]

% slanted
\margin{On suppose tjrs $M$ accessible dans $3^\circ$ et $4^\circ$.}
\textcircled{3} Pour qu'il existe un quadruplet qui,
avec $W$, définisse une pré-$Q$-catégorie de
modèles à droite $(TF=F\cap W)$\inserted{\textsuperscript{$*$}}, il faut et
% slanted
\margin{$*$ avec en plus $F,TF$ (et $C,TC$) closes à droite (resp à g.)}
il suffit qu'on ait b$'$) ci-dessus, ainsi que
\deleted{\unsure{la condition} \illegible{2}}
\[
 (W^{\mathrm{univ}})_*\cap W_{\mathrm{univ}}\ \text{\deleted{\unsure{est close}} \inserted{est clos à droite}}
\]
\deleted{(\unsure{cf. l'opération} $\unsure{\Phi|\underline{\varphi}|\underline{\mathcal{Q}}|\underline{\mathcal{D}}}\longmapsto\overline{\Phi}$), \unsure{terminologie|termologie} de la p.~1\unsure{*} et de la prop.~1 (p.~2).}
% archive sheet number: 20
\scanpage{22}
\mspage{16}

\textcircled{4} Les hypothèses préliminaires sur $(M,W)$, y
\inserted{\unsure{inclus|incluse} l'accessibilité,}
ainsi que les conditions équivalentes du \textcircled{2},
sont stables par passage de $(M,W)$ à
$(M(I),W(I))$. Mais qu'en est-il de la
condition du \textcircled{3} ? Est-elle conséquence
des conditions du \textcircled{2} (i.e. leur est-elle
équivalente ?).

Je vais esquisser la démonstration du
théorème\inserted{-scholie}, pour cerner le travail technique
qui reste à faire.

$1^\circ$) est de la vérification de routine.

$2^\circ$) Voici les implications $\pm$ tautologiques
entre les conditions envisagées :
\[
 \text{les}\qquad a''\Rightarrow a'\overset{?}{\Rightarrow}a\Rightarrow b\deleted{\unsure{\Rightarrow b'}}\Leftarrow b',
\]
mais je vois que $a'\Rightarrow a$ n'est \deleted{pas} nullement
``tautologique'', et peut-être même faux
après tout !\textsuperscript{$*$} Il faudrait supposer que dans
% slanted
\margin{\unsure{C'est pourquoi} j'ai modifié \unsure{$a'$)}, \unsure{en imposant|en renforçant} à a$'$) \deleted{\unsure{pour}} \unsure{l'égalité des|j'y ai mis} \unsure{?} après le trait !}
le quadruplet $(C,F,TC,TF)$, $F$ et $TF$ soient clos
à droite, \inserted{i.e. $\unsure{F=}\deleted{\unsure{TF}}$ et \uline{plus précisément} que $F=(TC)_*$, $TF=\unsure{,}C_*$,} ce qui sera le cas p.ex. si
on a une structure de $Q$-catégorie de
modèles ``close'' (``closed model category''), au
sens de Quillen. Les autres \inserted{trois} implications dans
\unsure{(*)} sont bel et bien tautologiques.

$a\Rightarrow a'$ On peut dans a) supposer $\deleted{\unsure{TC_0}}$
\[
 TC_0\subset C_0
\]
(quitte à remplacer $C_0$ par $C_0\cup TC_0$), puis, grâce à
% archive sheet number: 21
\scanpage{23}
\mspage{17}

$1^\circ$) ($W^{\mathrm{univ}}$ \unsure{et|$\cap$} $\mathrm{Cof}_W$ sont clos à gauche), \uline{posant}
\[
 \overline{TC_0}=TC,\quad \overline{C_0}=C,\quad \text{on aura}\quad
 \begin{aligned}TC&\subset W^{\mathrm{univ}},\\ C&\subset\mathrm{Cof}_W.\end{aligned}
\]

\deleted{(quitte à \unsure{barrer|remplacer} nos parties $TC_0,C_0$ par leurs
clôtures à gauche $TC=\overline{TC_0}$, $C=\overline{C_0}$), \unsure{Nous}
\unsure{pouvons supposer|peut suffire} qu'elles soient petites. \unsure{Dans le}
\illegible{3} de \unsure{la conclusion} \illegible{1}, \unsure{point},}

Passant de $TC_0$ à $TC$, de $C_0$ à $C$
on a perdu seulement que ces parties de
$\mathrm{Fl}\,M$ soient \unsure{petites}. Mais par le th.~1, posant
\[
 F=(TC)_*,\quad TF=C_*
\]
les \deleted{\unsure{deux}} couples de parties closes (à g. -- à \unsure{dr})
\[
 (C,TF),\quad (TC,F)
\]
satisfont à la condition de factorisation de
Quillen. On aura donc les propriétés
$M1$ à $M4$, plus les inclusions
$TC\subset C\cap W$ \unsure{,} (car $TC\subset W^{\mathrm{univ}}\subset\deleted{\unsure{W,}}W$ \inserted{et $TC\subset C$ par constr.}) et
$C\subset\mathrm{Cof}_W$ de $M'$ et $M''$. Il reste à vérifier
$TF\subset F\cap W$, $F\subset\mathrm{Fib}_W$, mais $F(=(TC)_*)\subset\mathrm{Fib}_W$ par hyp., et $TF(=\deleted{\unsure{W}}\unsure{C_*|C_W})\subset W$
aussi ; d'autre part $TF\subset F$ résulte
de $TC\subset C$.

$a'\Rightarrow a''$. On veut donc, en plus
des conditions $M1$ à $M4$ et $M',M''$, la
condition $TC=C\cap W$ $(=C\cap W^{\mathrm{univ}})$.
L'idée, c'est de remplacer le couple
% archive sheet number: 22
\scanpage{24}
\mspage{18}

\[
 C,TC\qquad(C\subset\mathrm{Cof}_W,\ TC\subset\deleted{\unsure{W^{\mathrm{univ}}\cap}}C\cap W)
\]
donnant naissance à
\[
 F=(TC)_*,\quad TF=C_*,
\]
par
\[
 C,\ C\cap W\quad(\supset TC)
\]
ce qui ne change pas le couple $(C,TF)$,
mais bien le couple $(TC,F)$, remplacé
par
\[
 \begin{aligned}
 (TC)' &={} C\cap W,& F'&=(TC')_*\\[-2pt]
 \cup &\deleted{\unsure{\cap|?|\,\uparrow}}\hspace{1.5em}\parallel &&\cap\\[-2pt]
 TC &\quad C\cap W^{\mathrm{univ}}&&F
 \end{aligned}
\]
% slanted
\margin{NB $TC'$ sera encore clos à gauche, comme intersection de deux parties closes à gauche.}
\unsure{--} \unsure{renforcé|remplacé} $F$ par un $F'$ plus petit).

Le seul problème, c'est si $\deleted{\unsure{(C,TF)}}$
$(TC',F')$ satisfait encore à la condition
de factorisation de Quillen. Ce sont des
parties closes (à gauche pour $TC'$, à dr.
pour $F'$) ``duales'' l'une de l'autre
\[
 TC'={F'}^{*},\quad F'=TC'_*,
\]
mais comme on ne sait pas \deleted{\unsure{d'avance|d'avan}}\inserted{si}
$TC'$ est \deleted{\unsure{engendré}} de la forme $\widetilde{TC'_0}$, avec
$TC'_0$ petite, on n'est pas sûr que la \deleted{\unsure{th.}}
conclusion du th.~1 s'applique. \deleted{\unsure{Mais}} \inserted{Mais}
si on suppose dans a$'$) que $C$ est
une partie \uline{accessible} de $\mathrm{Fl}\,M$ (i.e., étant
déjà close à g, que $C=\widetilde{C_0}$, $C_0$ petite),
% archive sheet number: 23
\scanpage{25}
\mspage{19}

alors je présume qu'on pourra prouver
qu'il en est de même de $C\cap W$ ($W$
elle-même étant accessible, grâce à la
condition de stabilité par petites $\varinjlim$
filtrantes), donc le th.~1 s'applique bel et
bien à $TC'=C\cap W$. Mais si on garde
\unsure{$a'$)|$a''$)} telle quelle, alors il faudrait sans doute
prouver que $a'$) \uline{implique} $a'$ sous la forme
plus forte, avec $C$ \uline{accessible}. \inserted{Mais} \unsure{cela on|cela car} on a $a'\Rightarrow a$ \deleted{\unsure{impl.}}, \deleted{il suffit pour cela}
\inserted{et qu'on vient d'établir} que $a\Rightarrow a'$ sous cette forme plus forte
\deleted{de \unsure{constater|montrer} que b) implique a$'$ sous cette forme
plus forte, ou encore que b implique a}
\unsure{de quoi} (modulo \unsure{le fait} \inserted{conjecturé} \unsure{ci-dessus} sur les
parties \unsure{admissibles|accessibles} de $\mathrm{Fl}\,M$).

Il reste donc à prouver
\[
 b\Rightarrow a,\qquad b\Rightarrow b',
\]
et il suffit de prouver la première implication,
car quand on sait que $a\Rightarrow a''$), et que $a'$)
implique que
\[
 \deleted{\unsure{\mathrm{Cof}_W\supset}}\underset{\unsure{\subset}}{\vphantom{F}}\mathrm{Fib}_W\cap W=W_{\mathrm{univ}}
\]
(cor. 1 page 10), donc $(\mathrm{Cof}_W)_*\subset W$ et $(W^{\mathrm{univ}})_*\subset\mathrm{Fib}_W$
\deleted{\unsure{impli}} (i.e. b)) impliquent, $(\mathrm{Cof}_W)_*\subset(W^{\mathrm{univ}})_*$
(car $W^{\mathrm{univ}}\subset\mathrm{Cof}_W$), \inserted{comme}
$(\mathrm{Cof}_W)_*\subset W\cap\mathrm{Fib}_W=W_{\mathrm{univ}}$,
ce qui est la partie à établir de b$'$) (l'autre
$(W^{\mathrm{univ}})_*\subset\mathrm{Fib}_W$ étant tautologique ---
\unsure{contenue|contenu} dans b)).
% archive sheet number: 24
\scanpage{26}
\mspage{20}

Notons maintenant que b) n'est autre que
la variante affaiblie de a), quand on omet
de spécifier que $C_0,TC_0$ doivent être \unsure{petites|petits}.
Donc le seul problème pour $b\Rightarrow a$,
c'est de voir ceci :

\uline{Lemme conjectural} \textcircled{1} \deleted{\illegible{2}}
\deleted{Si on a une partie
close à gauche $\unsure{\Phi|\underline{\varphi}|\underline{\mathcal{Q}}|\underline{\mathcal{D}}}$ de $\mathrm{Fl}\,M$, telle que
$\Phi_*\subset\mathrm{Fib}_W$ (resp. $\Phi_*\subset W$), alors il
existe une partie \uline{\unsure{petite}} $\Phi_0\subset\Phi$,
telle que $\Phi_{0*}\subset\mathrm{Fib}_W$ (resp. $\Phi_{\unsure{0|0*}}\subset W$).}
% slanted
\margin{Ce lemme a l'air un peu faux, mais ça doit pouvoir s'arranger\ldots}

\inserted{\textcircled{2} Si $C,C'$ sont deux parties closes à gauche de $\mathrm{Fl}\,M$ et
$C$ et $C'$ sont accessibles, il en est de \unsure{$=$} de $C\cap C'$, \unsure{Si $a''$)|il y a}}

Je crois que c'est \deleted{\unsure{bel}}
sous les hypothèses faites sur $M,W$ ---
qui doivent impliquer que dans un sens
convenable, $W$ est une partie ``accessible''
de l'ens $\mathrm{Fl}\,M$ des flèches de la cat.
accessible $M$. Il doit en résulter qu'une
partie de $M$ construite explicitement en
terme de $W$, telle $\mathrm{Fib}_W$, est elle-$=$
accessible, et le lemme doit sortir de
là.

Ainsi, $2^\circ$ est prouvé, modulo ce lemme
d'accessibilité, qu'il va falloir prouver
$\pm$ en $=$ \deleted{temps} que le th.~1 (que je n'ai
fait \unsure{ici} qu'énoncer) ;
% archive sheet number: 25
\scanpage{27}
\mspage{21}

Prouvons $3^\circ$. Supposons d'abord qu'il existe
un quadruplet ``clos'' $(C,TC,F,TF)$, s'insérant
dans un ``pré-quintuplet de Quillen clos'' qui
soit une \inserted{$Q$-}cat. de modèles à droite, i.e.
\[
 TF=F\cap W=F\cap W_{\mathrm{univ}}.
\]

\deleted{Comme} Prouvons donc
\[
 \deleted{F\subset\mathrm{Fib}_W,\ \text{on a}\quad TF=}
\]
\deleted{\illegible{3}}
\[
 \text{que}\qquad\underbrace{(W^{\mathrm{univ}})_*\cap W_{\mathrm{univ}}}_{A_W}
 \quad\text{\deleted{\unsure{est clos}} est clos à \unsure{droite}}
\]

En effet, on a
\[
 TC\subset W^{\mathrm{univ}}
\]
d'où
\[
 (W^{\mathrm{univ}})_*\subset(TC)_*=F
\]
donc
\[
 A=(W^{\mathrm{univ}})_*\cap W_{\mathrm{univ}}\subset F\cap W_{\mathrm{univ}}
 \underset{\substack{\uparrow\\\text{(par hyp.)}}}{=}TF,
\]
or $TF$ est clos à droite et contenu dans
$W_{\mathrm{univ}}$, d'où
\[
 \overline A\subset\deleted{\overline{\unsure{(W^{\mathrm{univ}})_*\cap W_{\mathrm{univ}}}}}\subset TF\subset W_{\mathrm{univ}},
\]
et comme $(W^{\mathrm{univ}})_*$ est clos à droite, on a aussi
\[
 \deleted{\overline{\unsure{(W^{\mathrm{univ}})_*\cap W_{\mathrm{univ}}}}}\overline A\subset(W^{\mathrm{univ}})_*,
\]
\deleted{q.e.d.}
% slanted
\inserted{d'où $\overline A\subset W_{\mathrm{univ}}\cap(W^{\mathrm{univ}})_*=A$, donc $A$ est clos à droite, qed.}

Prouvons la réciproque, i.e. supposons
satisfaites les conditions du $2^\circ$, \uline{plus} la
condition $(*)$ ci-dessus. On doit construire
\unsure{le|les} \unsure{qu. demandé|quadruplet}.

\deleted{On prendra
\[
 \begin{aligned}
 TC&=W^{\mathrm{univ}}&&\deleted{\unsure{TC}}\ \text{(le plus grand possible)}\\
 C&=\deleted{\unsure{\mathrm{Cof}_W}}\widetilde{(W_{\mathrm{univ}})^*\cup W^{\mathrm{univ}}}
 &&\text{(partie close}\\
 &&&\text{à g. engendrée par }(W_{\mathrm{univ}})^*\text{ et }W^{\mathrm{univ}}\text{)}
 \end{aligned}
\]
d'où $F=(TC)_*=(W^{\mathrm{univ}})_*$}
% archive sheet number: 26
\scanpage{28}
\mspage{22}

Notons d'abord, pour nous y reconnaître,
que les relations b$'$) du $2^\circ$, et les
relations transposées (compte tenu que $\mathrm{Cof}_W$
et $W^{\mathrm{univ}}$ sont closes à g., par $1^\circ$), \deleted{\unsure{donnent}}
savoir
\[
 \left\{\begin{aligned}
 (\mathrm{Cof}_W)_*&\subset W_{\mathrm{univ}},&
 (W^{\mathrm{univ}})_*&\subset\mathrm{Fib}_W\\
 (W_{\mathrm{univ}})^{\unsure{\alpha|*}}&\subset\mathrm{Cof}_W,&
 (\mathrm{Fib}_W)^{\alpha}&\subset W^{\mathrm{univ}}
 \end{aligned}\right.
\]
donnent\textsuperscript{$*$} les deux diagrammes \inserted{carrés} d'inclusions
% slanted
\margin{compte tenu de $W^{\mathrm{univ}}\subset\mathrm{Cof}_W$ $W_{\mathrm{univ}}\subset\mathrm{Fib}_W$ \unsure{$\}$} et des inclusions transposées}

\diagram{\begin{tikzcd}[ampersand replacement=\&,row sep=small,column sep=small]
 \& W^{\mathrm{univ}}\arrow[dr,hook] \& \\
 (\mathrm I)\qquad(\mathrm{Fib}_W)^*\arrow[ur,hook]\arrow[dr,hook] \& \& \mathrm{Cof}_W \\
 \& (W_{\mathrm{univ}})^{\unsure{\wedge|*}}\arrow[ur,hook] \&
\end{tikzcd}}

et
% slanted
\margin{\uline{NB} On a ces carrés d'inclusion, chaque fois qu'on a un syst. de Quillen $(W,C,F,TC,TF)$ \uline{clos}}

\diagram{\begin{tikzcd}[ampersand replacement=\&,row sep=small,column sep=small]
 \& (W^{\mathrm{univ}})_*\arrow[dl,hook,very thick] \& \\
 (\mathrm{II})\qquad\mathrm{Fib}_W \& \& (\mathrm{Cof}_W)_*\arrow[ul,hook]\arrow[dl,hook,very thick] \\
 \& W_{\mathrm{univ}}\arrow[ul,hook,"\alpha"'] \& 
\end{tikzcd}},

dont le premier est transposé (par $\unsure{\Phi|\underline{\varphi}|\underline{\mathcal{Q}}|\underline{\mathcal{D}}}\unsure{-|\to}\Phi^*$)
du second, et où les 8 \unsure{'}ensembles de
flèches envisagés, à l'exception de $\mathrm{Fib}_W$ et
$W_{\mathrm{univ}}$ (reliés par $\alpha$ dans le carré $\mathrm{II}$)
sont clos. (\deleted{\unsure{les} \unsure{de} \unsure{parties}} \inserted{les quatre} ensembles de
$\mathrm I$ clos à g., et les deux ens. de $\mathrm{II}$
non \unsure{intéressés} par $\alpha$, clos à dr.).
% archive sheet number: 27
\scanpage{29}
\mspage{23}

Les ``pré quintuplets de Quillen clos'', associés
à $W$, correspondent au choix de
parties closes à gauche $TC$ et $C$,
s'insérant dans le diagramme (I) comme
montré \inserted{en $\mathrm I'$} ci-dessous, et donnant lieu
par transposition à $F$ et $TF$ comme dans
$\mathrm{II}'$ ci-dessous

\diagram{\begin{tikzcd}[ampersand replacement=\&,row sep=small,column sep=small]
 \& \& W^{\mathrm{univ}}\arrow[ddrr,hook] \& \& \\
 \& TC\arrow[ur,hook]\arrow[ddrr,hook] \& \& \& \\
 (\mathrm I')\qquad(\mathrm{Fib}_W)^{\unsure{v|*}}\arrow[ur,hook]\arrow[ddrr,hook] \& \& \& \& \mathrm{Cof}_W \\
 \& \& \& C\arrow[ur,hook] \& \\
 \& \& (W_{\mathrm{univ}})^{\unsure{*|k}}\arrow[ur,hook,no head] \& \&
\end{tikzcd}}

\diagram{\begin{tikzcd}[ampersand replacement=\&,row sep=small,column sep=small]
 \& \& (W^{\mathrm{univ}})_*\arrow[dl,hook] \& \& \\
 \& F\arrow[dl,hook] \& \& \& \\
 (\mathrm{II}')\qquad\mathrm{Fib}_W \& \& \& \& (\mathrm{Cof}_W)_{\unsure{\alpha|*}}\arrow[uull,hook]\arrow[dl,hook,thick] \\
 \& \& \& TF\arrow[uull,hook]\arrow[dl,hook,thick] \& \\
 \& \& W_{\mathrm{univ}}\arrow[uull,hook] \& \&
\end{tikzcd}}\unsure{;}

le ``\unsure{hic}'' étant donc de choisir $TC$ et $C$
assez grandes pour qu'on ait
\[
 \begin{aligned}
 (TC)_*&\subset\mathrm{Fib}_W,&(\deleted{\unsure{TF}}C)_*&\subset W_{\mathrm{univ}}\quad-\\
 \parallel &&\parallel\\
 F&&TF
 \end{aligned}
\]
% archive sheet number: 28
\scanpage{30}
\mspage{24}

et les \deleted{\unsure{cond.}} inclusions stipulées dans $2^\circ$, b$'$)
assurent qu'on peut bel et bien trouver
$TC$ et $C$, quitte à les prendre les
plus grandes possibles $TC=W^{\mathrm{univ}}$, $C=\mathrm{Cof}_W$
--- auquel cas, d'ailleurs, on aura
\[
 TC=C\cap W^{\mathrm{univ}}_{\unsure{-}},
\]
i.e. on a une \inserted{pré-$Q$-}catégorie de \deleted{\unsure{Qu}}\unsure{modèles} à
g. Mais il s'agit d'assurer maintenant,
par contre, que
\[
 TF=F\cap W_{\mathrm{univ}}\unsure{)|3)|?},
\]
\deleted{Prenons}
ce qui implique, on vient de le voir \unsure{$*$} (en
notant $(W^{\mathrm{univ}})_*\subset F$) que
\[
 \begin{aligned}
 \overline{(W^{\mathrm{univ}})_*\cap W_{\mathrm{univ}}}
 &\subset\overline{F\cap W_{\mathrm{univ}}}=\overline{TF}={}&TF\\[-2pt]
 &&\cap\\[-2pt]&&W_{\mathrm{univ}}.
 \end{aligned}\tag*{($*$)}
\]

Inversement, si $(*)$ est satisfait, prenons
\[
 \begin{aligned}
 F&=(W^{\mathrm{univ}})_*\quad\text{le plus petit possible}\\
 &\text{i.e.}\quad C=W^{\mathrm{univ}}\quad\text{le plus \unsure{grd|gnd} poss.}
 \end{aligned}
\]
On devra donc prendre
\[
 TF=F\cap W_{\mathrm{univ}}=(W^{\mathrm{univ}})_*\cap W_{\mathrm{univ}},
\]
\deleted{\unsure{alors}} \inserted{et} cet ensemble est bien clos à droite \unsure{?}
\deleted{\unsure{pour} \illegible{4}} \inserted{par hypothèse.}
\deleted{\unsure{pour} \illegible{2}}

\uline{Corollaire} \inserted{Supposons les conditions b$'$) du th. Scholie \unsure{satisfaites}} Les conditions suivantes \inserted{a), b) sur $W$} sont
équivalentes (et impliquent \unsure{la|les} cond. du $3^\circ$ du th.-Scholie)
\deleted{a) On a \unsure{(et} aussi celles du $2^\circ$) du th. Scholie}
% archive sheet number: 29
\scanpage{31}
\mspage{25}

\textcircled{a}
\[
 \deleted{(W^{\mathrm{univ}})_*\subset\mathrm{Fib}_W}\ \text{et}\quad
 (W^{\mathrm{univ}})_*\cap\unsure{\overline{W}|W}^{\deleted{\unsure{\mathrm{univ}}}}_{\mathrm{univ}}=(\mathrm{Cof}_{\unsure{W|M}})_*
\]

\textcircled{b}
\[
 \deleted{\left\{\begin{aligned}
 (W^{\mathrm{univ}})_*&\subset\mathrm{Fib}_W\ \text{et}\\
 (\mathrm{Cof}_W)_*&\subset W_{\mathrm{univ}}
 \end{aligned}\right.}\qquad
 \underbrace{\widetilde{W^{\mathrm{univ}}\cup(W_{\mathrm{univ}})^*}=\mathrm{Cof}_W}
 _{\substack{\text{i.e. }\mathrm{Cof}_W\text{ est }\deleted{\text{\unsure{engendré}}}\inserted{\text{la partie}}\\
 \text{close à g. engendrée par}\\W^{\mathrm{univ}}\ \text{et}\ (W_{\mathrm{univ}})^*}},
\]

et ces conditions impliquent \inserted{\textcircled{c}} qu'il existe
une structure de \uline{triple \deleted{\illegible{1}}\inserted{clos} de Quillen}
$(\deleted{\unsure{F}}W,C,F)$ induisant $W$ (donc avec
$TC=W\cap C$, $TF=W\cap F$) (lequel sera clos au
sens de Quillen ssi $W$ stable par
facteurs directs --- car $C$ et $F$ sont déjà
stables spontanément).

Il suffit en effet de prendre donc
\[
 TC=W^{\mathrm{univ}},\quad C=\mathrm{Cof}_W.
\]
Mais je doute que \unsure{-} l'existence d'un
triple clos de Quillen $(W,C,F)$ suffise
à impliquer les conditions a), b)\unsure{$\ne$|,|?}
ci-dessus, ni $=$ la condition plus faible
\[
 (W^{\mathrm{univ}})_*\cap W_{\mathrm{univ}}=(\mathrm{Cof}_W)_*
\]
qui \deleted{suffit déjà} signifie \uline{exactement} que
$C=\mathrm{Cof}_W$, $TC=W_{\deleted{\unsure{\mathrm{univ}}}}^{\inserted{\mathrm{univ}}}$ et leurs transposées
$TF=(\mathrm{Cof}_W)_*$, $F=(W^{\mathrm{univ}})_*$ forment une \deleted{\unsure{syst.}} catégorie
de Quillen close
% archive sheet number: 30
\scanpage{32}
\mspage{\textcircled{26}}

Conditions sur $(M,W)$ avec \deleted{\unsure{des}}

\deleted{\unsure{(1)}} a) $M$ stable par $\varprojlim$ finies et $\varinjlim$ (petites).

b) $W$ \deleted{\unsure{il en est}} stable par $\varinjlim$ filtrantes

c) $M$ accessible.

\diagram{\scriptsize\begin{tikzcd}[ampersand replacement=\&,row sep=small,/tikz/scale=.82,/tikz/transform shape]
 {\mathrm I\left\{\begin{aligned}
 &\text{\textcircled{a} existence d'un syst. de Quillen faible }(W,C,TC,F,TF)\ \text{(clos si on y tient)}\\
 &\Updownarrow\\
 &\text{\textcircled{b} existence d'un syst. de Quillen faible \unsure{ex.ct|ex.t|exact} à gauche }(TC=C\cap W)\\
 &\Updownarrow\\
 &\text{\textcircled{c} }\underbrace{(\mathrm{Cof}_W)_*\subset W,\ \deleted{\unsure{\mathrm{Cof}}}(W^{\mathrm{univ}})_*\subset\mathrm{Fib}_W}_{(*)}
 \ \Longleftrightarrow\ (\mathrm{Cof}_W)_*\subset W_{\mathrm{univ}},\ (W^{\mathrm{univ}})_*\subset\mathrm{Fib}_W
 \end{aligned}\right.}\\
 {\mathrm{II}\left\{\begin{aligned}
 &\text{\textcircled{a} existence d'un syst. de Quillen faible \unsure{ex.t|ex.ct|exact} à droite }(TF=F\cap W)\\
 &\Updownarrow\\
 &\text{\textcircled{\unsure{d|b}} On a }(*),\ \text{et de plus }(W^{\mathrm{univ}})_*\cap W_{\mathrm{univ}}\ \text{clos à droite}\\
 &\Updownarrow\\
 &\deleted{\unsure{\mathrm{III}}}\ \text{\textcircled{b} existence d'un syst. de Quillen clos (closed model category) :}\\
 &\qquad TC=C\cap W,\ TF=F\cap W,\ F=(TC)_*,\ TF=C_*,\ TC=F^*,\ C=(TF)^*\\
 &\tikz[baseline=-.5ex]{\draw[<->,double] (-.19,.13)--(.19,-.13);}\\
 &\text{\textcircled{c} Posant }TC=W^{\mathrm{univ}},\ F=(TC)_*,\ TF=F\cap W_{\mathrm{univ}},\ C=(TF)_*,\\
 &\qquad\text{on trouve avec }W\text{ un syst. de Quillen clos}
 \end{aligned}\right.}\arrow[u,Rightarrow]\\
 {\mathrm{III}\left\{\begin{aligned}
 &C=\mathrm{Cof}_W\ \text{et }TC=W^{\mathrm{univ}}\ \text{font partie \inserted{(avec $W$)} d'un syst. de Quillen clos}\\
 &\qquad(\text{donc }F=(W^{\mathrm{univ}})_*,\ TF=(\mathrm{Cof}_W)_*,\ W=TF\circ TC\unsure{)}\\
 &\Updownarrow\\
 &(*),\ \text{et de plus }(W^{\mathrm{univ}})_*\cap W_{\mathrm{univ}}=(\mathrm{Cof}_W)_*
 \end{aligned}\right.}\arrow[u,Rightarrow]\\
 {\mathrm{IV}\left\{\begin{aligned}
 &\text{On a }(*)\text{ et }\mathrm{Cof}_W=\widetilde{(W^{\deleted{\unsure{\mathrm{univ}}}}_{\mathrm{univ}})^*\cup W_{\deleted{\unsure{\mathrm{univ}}}}^{\mathrm{univ}}}\\
 &\Updownarrow\\
 &\text{On a }(*)\text{ et }(W^{\mathrm{univ}})_*\cap\overline{W_{\mathrm{univ}}}=(\mathrm{Cof}_W)_*
 \end{aligned}\right.}\arrow[u,Rightarrow]
\end{tikzcd}}

Cas $M=\mathrm{Cat}$, $W=W_\infty$. On sait que \unsure{II|III} est satisfait, par
Thomason. Donc on a \deleted{\unsure{III} \illegible{1}}
\[
 (\mathrm{Cof}_W)_*\subset W_{\mathrm{univ}},\quad(W^{\mathrm{univ}})_*\subset\mathrm{Fib}_W,\quad
 (W^{\mathrm{univ}})_*\cap W_{\mathrm{univ}}\ \text{clos à \unsure{droite}}
\]
J'ai pu prouver directement les deux premières inclusions
(i.e. $\mathrm I$), mais pas que $(W^{\mathrm{univ}})_*\cap W_{\mathrm{univ}}$ soit clos
à droite. (Je devrais au moins essayer !)

\deleted{\unsure{Je} \unsure{signale que} \illegible{2} \unsure{du tout} \unsure{on a}} \inserted{On n'a \uuline{pas} (IV), i.e. on a}
\[
 (W^{\mathrm{univ}})_*\cap\overline{W_{\mathrm{univ}}}\ne(\mathrm{Cof}_W)_*
 \deleted{\unsure{(W^{\mathrm{univ}})_*\cap\overline{W_{\mathrm{univ}}}}}
\]
car
\[
 (W^{\mathrm{univ}})_*\cap\overline{W_{\mathrm{univ}}}\not\subset W_{\mathrm{univ}}
\]
(contient les $X\to e$ avec $\unsure{\alpha|X}$ discrète)
donc on a
\[
 (W^{\mathrm{univ}})_*\cap\overline{W_{\mathrm{univ}}}\supsetneq
 \underbrace{(W^{\mathrm{univ}})_*\cap W_{\mathrm{univ}}}_{\substack{\text{\unsure{n'est pas l'orth.} de}\\W^{\mathrm{univ}}\cup(W_{\mathrm{univ}})^*}}
\]
\margin{Il suffirait de savoir que $\overline{W_{\mathrm{univ}}}\cap\mathrm{Fib}_W=W_{\mathrm{univ}}$ ce qui impliquerait que III et IV sont équivalents}
\inserted{Mais c'est faux, p.ex. si $X$ cat. discr. non vide, $X\to e$ est dans $\overline{W_{\mathrm{univ}}}\cap\mathrm{Fib}_W$}
% archive sheet number: 31
\scanpage{33}
\mspage{71}

excellence !). Aussi est-il rare que la raison n'en soit elle-même profondément pervertie, et elle l'est en tous cas toujours peu ou prou.

Il n'y a pas à s'en étonner, quand on se rappelle l'omniprésence des pulsions de désirs dans la psyché, et notamment dans toutes les activités créatrices, tant chez soi-même que chez autrui. Aussi, quand la relation au sexe est faussée, toute la vision de la réalité, et avec elle toutes les relations humaines, sont faussées. Ainsi la raison se trouve sommée, soit de démissionner purement et simplement (c'est là souvent l'attitude du mystique), soit de ``rationaliser'' (!) une vision et des relations aberrantes, par des arguments et des considérations tout aussi aberrants. Même chez des êtres d'une intelligence exquise, quand telle pieuse fiction est en jeu, soudain on se trouve confronté à une débilité telle qu'on croit rêver...

35. La Fuite est le remède - ou la Maison des fous

La double-pensée est un exemple typique et extrême de ``\textit{la Fuite}'' dans la psyché. Nous entendons par là un ensemble cohérent de mécanismes, d'attitudes et de comportements de la psyché, qui inexorablement et de façon permanente nous poussent à nous affirmer (avec une force véhémente dans les couches inconscientes, et comme chose allant de soi dans le conscient) des choses qui sont soit manifestement absurdes ou dénuées en elles-mêmes de tout sens (mais dont le \textit{vrai} sens réside dans l'intention inconsciente de brouiller, de noyer un poisson, ou de faire taire une inquiétude...), soit encore à l'exact opposé de \unsure{.|:}ce qu'au fond nous savons pertinemment, ou qui en représentent une déformation flagrante. (Et de plus dans un aspect si essentiel, le plus souvent, que cette version de la réalité que nous nous présentons à nous-même n'en représente plus qu'une carricature plus ou moins grotesque). \textit{Cet étrange syndrome de la Fuite domaine silencieusement et inexorablement la vie de quasiment tous les hommes, de façon plus ou moins tyrannique et totale suivant le degré de maturité et l'état de fidélité ou d'infidélité de l'âme} - et pourtant, d'un commun et universel accord, il est ignoré par tous ! C'est dire que nous fonctionnons et évoluons, avec les airs désinvoltes de gens très distingué et qui savent de quoi ils parlent, dans une gigantesque \textit{Maison de Fous} à l'échelle de la planète entière, sans autres gardes-chiourme que la commune Folie de tous, et où tous les humains seraient les aliénés et les fous-à-lier et internés à vie.

Quand l'âme pour la première fois s'éveille à la réalité de la Maison de Fous, de cette \textit{Fuite} universelle, obstinée, de tous les instants devant tout ce qui effraie la psyché - elle en reste saisie, comme pétrifiée. Il faudra un bon moment avant qu'elle ait fini de se tater et qu'elle en croie ses yeux et ses saines facultés ! Car elle est seule à voir ce qu'elle voit, et elle ne peut s'empêcher de se dire que c'est elle sûrement qui doit être toquée et que tous ces gens si posés et si sérieux et tous d'accord, au fond, sont sains d'esprit.

Dans l'histoire des hommes sur Terre, le premier à oser en croire le témoignage de ses yeux, en voyant la Maison de Fous, a été \textit{Freud}, au moment où il a découvert l'inconscient. Cela a été (me confirme Mère) la plus grande découverte de l'esprit humain sur terre. Cela a été aussi un des actes essentiels pour marquer le commencement du Temps \deleted{de} \unsure{.}la Fin, et pour préparer le Nouvel Age - l'âge de la Libération. C'est cet \unsure{.}acte, et l'oeuvre de fidélité qui en est née, qui ont joué le rôle le plus essentiel dans l'Oeuvre de Salut et de Libération préparée par la mission du Bouddha et
% archive sheet number: 32
